\section{Discriminant \texorpdfstring{\(-8\)}{-8}}\label{sec:minus8}

Put
\[
 Q_8(x,y)=x^2+2y^2,\qquad
 P_8=\{Q_8(x,y)>0:\gcd(x,y)=1\},
\]
and let \(C_8\) contain \(1\) and the positive integers whose prime
factors are \(1\) or \(3\pmod8\).

\begin{theorem}\label{thm:minus8-exact}
There are no exceptional admissible targets:
\[
\begin{aligned}
 C_8+C_8&=\{N\ge2:N\ {\rm even},\ 8\nmid N\},\\
 C_8+2C_8&=\{N\ge3:N\ {\rm odd}\},\\
 P_8+P_8&=\{N\ge2:16\nmid N\}.
\end{aligned}
\]
For every \(N\ge2^{78}\) with \(16\nmid N\),
\[
              \#\{(a,b)\in P_8^2:a+b=N\}
                       \ge N^{2/3}/4800 .
\]
\end{theorem}

Write \(Q,P,C\) for these objects in this section.
The exact dyadic shell formulas will supply both a large-core
theorem and an analytic prime boundary. Sharper selected-shell
bounds then reduce the remaining work to \(107\) rational interval
tests and complete finite core coverage.

\subsection{Primitive values and dyadic reduction}

The ring \(R=\Z[\sqrt{-2}]\) is norm Euclidean:
rounding both coordinates leaves error norm at most
\(1/4+2/4=3/4<1\). Thus it is principal.
An inert prime, congruent to \(5,7\pmod8\), dividing a norm
divides both coordinates. Primitivity permits the ramified
prime ideal of norm two to exponent at most one.
For split primes, orient each full prime power to one of the two
conjugate principal prime ideals. These observations prove
\begin{equation}
                        P=C\mathbin{\dot\cup}2C.               \label{eq:m8-value}
\end{equation}
Every \(C\)-value is \(1\) or \(3\pmod8\). Consequently
\begin{equation}
                   P+P=(C+C)\cup2(C+C)\cup(C+2C).              \label{eq:m8-union}
\end{equation}
The first two terms have residues respectively \(2,4,6\pmod8\)
and \(4,8,12\pmod {16}\), whereas the last has odd targets
at least three. This proves every local exclusion in the theorem;
the two core identities imply the full identity.

Multiplication by \(\varpi=\sqrt{-2}\) acts by
\[
                   (x,y)\longmapsto(-2y,x).
\]
It bijects primitive odd-norm vectors with primitive even-norm
vectors; the latter have even first coordinate and odd second
coordinate, and the inverse is \((X,Y)\mapsto(Y,-X/2)\).
At odd primes the determinant is a unit. Thus every division or
multiplication by two below is justified on the indicated parity
class; no general closure of primitive vectors under doubling is
being assumed.

\subsection{Exact dyadic shell formulas}

Let \(r(N)\) count the full signed ordered \(Q\oplus Q\) shell,
let \(s(N)\) count \(Q\oplus4Q\), and let \(G_2(N)\) count shell
points for which neither binary coordinate pair is divisible by
two. Only dyadic primitivity is imposed in \(G_2\).
With \(\sigma(t)=0\) for nonpositive or nonintegral \(t\), define
\[
 f_{32}(\tau)=\eta(4\tau)^2\eta(8\tau)^2
     =\sum_{N\ge1}a_Nq^N
     =q\prod_{j\ge1}(1-q^{4j})^2(1-q^{8j})^2 .
\]
Then
\begin{equation}
\begin{aligned}
 r(N)&=4\sigma(N)-4\sigma(N/2)+8\sigma(N/4)-32\sigma(N/8),\\
 s(N)&=\sigma(N)-\sigma(N/2)+8\sigma(N/16)-32\sigma(N/32)+a_N,\\
 G_2(N)&=r(N)-2s(N)+r(N/4).                                  \end{aligned}\label{eq:m8-theta}
\end{equation}
We include the modular verification, so these formulas do not
depend on extrapolating a finite list of shell counts.

The even Gram matrices of \(Q\oplus Q\) and \(Q\oplus4Q\)
have levels \(8,32\) and square determinants \(64,1024\).
Their theta series therefore belong respectively to
\(M_2(\Gamma_0(8))\) and \(M_2(\Gamma_0(32))\), with trivial
character. Put
\(E_2(\tau)=1-24\sum_{n\ge1}\sigma(n)q^n\).
The proposed Eisenstein series are
\[
\begin{aligned}
 E_r&=-\tfrac1{24}
       \{4E_2(\tau)-4E_2(2\tau)+8E_2(4\tau)-32E_2(8\tau)\},\\
 E_s&=-\tfrac1{24}
       \{E_2(\tau)-E_2(2\tau)+8E_2(16\tau)-32E_2(32\tau)\}.
\end{aligned}
\]
To verify holomorphy at all cusps, write
\[
 D_d(\tau)=dE_2(d\tau)-E_2(\tau)
    =\frac1{2\pi i}\frac d{d\tau}
           \log\frac{\Delta(d\tau)}{\Delta(\tau)} .
\]
The quotient is a modular function, nonzero in the upper
half-plane and meromorphic at cusps. If at a cusp of width \(h\)
it is \(q_c^\nu U(q_c)\), \(U(0)\ne0\), its logarithmic derivative
has expansion
\[
                    \frac1h\left(\nu+
                              q_c\frac{U'(q_c)}{U(q_c)}\right).
\]
There is no cusp pole. Hence \(D_d\) is holomorphic of weight two,
and
\[
 E_r=D_2/12-D_4/12+D_8/6,\qquad
 E_s=D_2/48-D_{16}/48+D_{32}/24
\]
have the required levels. Both have constant term one.

Martin--Ono \cite[Theorem 1, level-32 row]{MartinOno} identify
\(f_{32}\) as a normalized weight-two newform.
The Sturm bounds are \(2\) at level \(8\) and \(8\) at level
\(32\) \cite[Theorem 1]{Sturm}. The necessary coefficient lists
are
\[
\begin{array}{c|rrrrrrrrr}
 N&0&1&2&3&4&5&6&7&8\\ \hline
 {}[q^N]\Theta_{Q\oplus Q}&1&4&8&&&&&&\\
 {}[q^N]E_r&1&4&8&&&&&&\\
 {}[q^N]\Theta_{Q\oplus4Q}&1&2&2&4&4&4&8&8&8\\
 {}[q^N](E_s+f_{32})&1&2&2&4&4&4&8&8&8
\end{array}
\]
They follow by direct lattice enumeration, divisor sums, and
truncated multiplication of the displayed eta product.
All coefficients are integers. Apply Sturm's congruence theorem
modulo every rational prime: each coefficient of the difference
is divisible by every prime, so is zero.
This proves the first two formulas in \eqref{eq:m8-theta};
the last is inclusion-exclusion.

The CM formula in \cite[Theorems 2 and 3]{MartinOno} gives
\(|a_p|\le2\sqrt p\) at every odd prime, with \(a_p=0\) for
\(p\equiv3\pmod4\). Its unramified Hecke roots have modulus
\(\sqrt p\); the Euler recurrence and multiplicativity imply
\(|a_N|\le\tau(N)\sqrt N\) for odd \(N\).
The eta product is supported on \(N\equiv1\pmod4\), so its
even coefficients vanish.

Write \(N=2^em\), \(m\) odd. For \(0\le e\le3\),
\[
 r(N)=4c_e\sigma(m),\qquad (c_0,c_1,c_2,c_3)=(1,2,6,6),
\]
and
\begin{equation}
 (G_2(m),G_2(2m),G_2(4m),G_2(8m))
   =(2\sigma(m)-2a_m,\ 4\sigma(m),\
                               20\sigma(m),\ 16\sigma(m)).     \label{eq:m8-dyadic}
\end{equation}
For \(e\ge4\), \(G_2(N)=0\).

\subsection{The nonmaximal dyadic order and Hecke bound}

Let \(H\) be the Hurwitz order in Hamilton's quaternion algebra,
with generators \(I,J,K\). Put \(\iota=I+J\), \(j=K\), and
\[
 O=\Z[\iota]\oplus\Z[\iota]j
   =\{a+bI+cJ+dK:a,b,c,d\in\Z,\ b\equiv c\pmod2\}.
\]
Here \(\iota^2=-2\), \(j^2=-1\), \(j\iota=-\iota j\),
and the norm is \(Q\oplus Q\).
This order has index four in \(H\); it is not maximal at two.

We verify its locally principal right-ideal class number directly.
The Hurwitz order is norm Euclidean. After rounding a real
quaternion to an integer vector, let its coordinate errors be
\(|r_i|\le1/2\), and put \(s=\sum r_i^2\).
The squared distance to an appropriate half-integer vector is
\(s-\sum|r_i|+1\). Since \(s\le\frac12\sum|r_i|\), one of the
two distances is at most \(1/2\).
Thus all Hurwitz right ideals are principal.

Away from two, \(O_\ell=H_\ell\).
Normalize the additive volume of \(H_2\) to one.
Its half-integer coset consists of odd-norm elements, as does
half its integer coset; hence
\(\operatorname{vol}(H_2^\times)=3/4\).
The additive volume of \(O_2\) is \(1/4\), and half its elements
have odd norm, so
\(\operatorname{vol}(O_2^\times)=1/8\).
The unit index is six. The \(24\) Hurwitz units have stabilizer
\[
                    O^\times=\{\pm1,\pm K\}
\]
of size four on these cosets, so act transitively.
Principalization in \(H\), followed by this unit transitivity
at two, principalizes every locally principal \(O\)-ideal.

The dyadic unit norm map is onto \(\Z_2^\times\).
Indeed \(1\), \(1+\iota\), \(1+\iota+\iota j\),
\(1+\iota+2j\) have norms \(1,3,5,7\).
Their inverses belong to \(O_2\).
Every \(2\)-adic unit is one of these four norms times a
unit square, supplied by a rational \(2\)-adic scalar.

For odd inert \(p\), reduction gives \(O/pO=M_2(\mathbb F_p)\).
The projective unit image \(\Gamma_p\) has order two, and the
fine quotient is \(Y_p=\Gamma_p\backslash\mathrm{PGL}_2(\mathbb F_p)\),
with uniform measure. Class number one identifies this quotient
with the adelic congruence quotient: its stabilizer is precisely
the projective unit image. To check this, local scalar
corrections have even norm valuations and combine into one
rational scalar, leaving an \(O\)-unit.
For \(p\nmid n\), define concretely
\begin{equation}
 (T_nh)(\Gamma_pg)
    =\frac14\sum_{\operatorname{nrd}\alpha=n}
                             h(\Gamma_p\rho_p(\alpha)g).       \label{eq:m8-Hecke}
\end{equation}
For odd \(m\), this is the all-integral-ideal operator, of degree
\(\sigma(m)\). Inverse principal generators give the displayed
orientation by quaternion conjugation. The global norm characters
on \(Y_p\) are exactly \(1\) and the determinant square character
\(\chi_p\). The norm surjectivity just proved prevents an
additional dyadic character. At \(p\), central scalars and the
reduction kernel leave only determinant square classes.
Lemma~\ref{lem:hecke-transfer}, including its strong-approximation
exclusion of finite-dimensional split-place factors, therefore yields
\begin{equation}
             \|T_m\|_{\{1,\chi_p\}^{\perp}}
                                  \le\tau(m)\sqrt m
                  \quad(m\ {\rm odd},\ p\nmid m).             \label{eq:m8-Ramanujan}
\end{equation}
All primes dividing \(m\) have maximal split local level.
No unramified theorem at two has been used.

For \(0\le e\le3\), let \(A_e=T_{2^e}\) be given by
\eqref{eq:m8-Hecke}. It is self-adjoint with nonnegative entries,
row sum \(c_e\), and norm at most \(c_e\).
It has \(\chi_p\)-eigenvalue
\((2/p)^ec_e\). Coprime local factorization gives
\begin{equation}
                       T_{2^em}=A_eT_m=T_mA_e.                \label{eq:m8-factor}
\end{equation}
For example, for \(\alpha\in O\) of norm \(2^em\), the right
superideal with localization \(O_2\) at two and
\(\alpha O_\ell\) at odd primes is integral, locally principal,
and of norm \(m\). Write it as \(\beta O\),
\(\operatorname{nrd}\beta=m\); then
\(\alpha=\beta\gamma\), \(\operatorname{nrd}\gamma=2^e\).
Exactly four choices differ by an \(O\)-unit. Reversing the
coprime local factors gives commutation. Thus the complementary
norm of \(T_N\) is at most \(c_e\tau(m)\sqrt m\).

\subsection{The uniform sieve and its infinite ranges}

Put
\[
\begin{aligned}
 \mathcal I&=\{p\ {\rm prime}:p\equiv5,7\pmod8\},\\
 E(m)&=\frac{\tau(m)\sqrt m}{\sigma(m)},\\
 M(X)&=\max_{1\le t\le X}
          \prod_{\substack{\ell\mid t\\\ell\equiv1,3(8)}}
                                      (v_\ell(t)+1).
\end{aligned}
\]
Let \(b_p(N)\) count full-shell points with at least one binary
coordinate pair divisible by \(p\).
For \(p\nmid N\), that condition is the nonsplit-torus normalizer
\(H_p\), of order \(2(p+1)\).
Its fraction in \(\mathrm{PGL}_2(\mathbb F_p)\) is
\(\mu_p=2/(p(p-1))\).
Its indicator \(u\) has mean \(\mu_p\), zero coefficient on
\(\chi_p\), and
\(\|u-\mu_p\|^2=\mu_p(1-\mu_p)\), since the torus norm is onto.
The operator commutes with the right action and is constant on
the support of \(u\). Thus
\(\langle u,T_Nu\rangle=\mu_pb_p(N)/4\).
Equations \eqref{eq:m8-Ramanujan}--\eqref{eq:m8-factor} give
\begin{equation}
                  \frac{b_p(N)}{r(N)}
                      \le\mu_p+(1-\mu_p)E(m)
                      \le\mu_p+E(m).                          \label{eq:m8-small}
\end{equation}
If \(p\mid N\), anisotropy forces both blocks divisible by \(p\).
The spoiled count is zero at valuation one; otherwise its ratio
is \(\sigma(m/p^2)/\sigma(m)<p^{-2}\le\mu_p\).
This treats every target valuation without applying a
level-\(p\) estimate to an index divisible by \(p\).

For a first bound on large primes,
\[
  \#\{(x,y):0<x^2+2y^2\le X\}
       \le\pi X/\sqrt2+(2+\sqrt2)\sqrt X\le4X .
\]
The last inequality holds for \(X\ge4\); the three smaller
integer bounds give \(2,4,8\) points. Flooring gives real \(X\).
In the \(G_2\)-shell both blocks are nonzero, and
\(r_Q(t)\le2M(N)\). Choose the divisible block, count its
quotient vector, and use this bound for the other block.
The large-\(p\) spoiled proportion is at most
\begin{equation}
                    L_eM(N)p^{-2},\qquad
               L_e=\frac{4\cdot2^e}{c_e}\le\frac{16}{3}.       \label{eq:m8-coarsetail}
\end{equation}
Also
\begin{equation}
\begin{aligned}
 B:=\sum_{p\in\mathcal I}\mu_p&<3/16,\\
 \Pi(P):=\#\{p\in\mathcal I:p\le P\}&\le P/4+1,\\
 \sum_{\substack{p\in\mathcal I\\p>P}}p^{-2}
                    &\le1/(4P)+2/P^2.                       \end{aligned}\label{eq:m8-progression}
\end{equation}
The last two statements follow from the progressions \(5,7\pmod8\).
For the first, sum inert primes below \(64\) exactly and bound
the rest by
\[
 \frac52\left(\frac1{69^2}+\frac1{8\cdot69}
                  +\frac1{71^2}+\frac1{8\cdot71}\right).
\]
The resulting bound is
\(49287299505759574293/269836951018231737140<3/16\).

By \eqref{eq:m8-dyadic}, \(G_2(N)/r(N)\ge1/2-E(m)/2\).
A sufficient condition for a positive surviving shell is
\begin{equation}
 B+E(m)(P/4+3/2)
      +L_eM(N)\{1/(4P)+2/P^2\}<1/2.                          \label{eq:m8-coarse}
\end{equation}
Surviving norms are positive, have no inert factor, and satisfy
the dyadic condition in \eqref{eq:m8-value}; their values lie in \(P\).

The complete finite Euler products give
\begin{equation}
\begin{aligned}
 M(X)&\le6X^{1/6},& M(X)&\le2X^{1/4},\\
 \tau(m)&\le7m^{1/6}\quad(m\ {\rm inert\!-\!supported}),&
 \tau(m)&\le42m^{1/6}\quad(m\ {\rm odd}).                       \end{aligned}\label{eq:m8-Euler}
\end{equation}
For the sixth-power bounds, maximize
\((a+1)^6/p^a\) at each odd prime below \(64\), over \(0\le a\le6\).
The factors decrease afterwards, and factors at primes at least
\(64\) are at most one by \(a+1\le2^a\).
The split and inert products are respectively below \(6^6,7^6\);
their product gives \(42^6\).
For the fourth-power split bound only primes \(3,11\) contribute,
and their product is \(4096/297<2^4\).
These calculations establish \eqref{eq:m8-Euler} for all \(X,m\).

Since \(N\le8m\), \(M(N)\le9m^{1/6}\).
For arbitrary odd \(m\), use \(P=m^{1/4}\) in
\eqref{eq:m8-coarse}. The terms after \(B\) are at most
\[
                      \frac{45}{2}m^{-1/12}+159m^{-1/3}.
\]
If \(m\ge2^{75}\), use \(m^{1/12}>224/3\) to obtain
\begin{equation}
                    \frac3{16}+\frac{135}{448}
                                    +\frac{159}{2^{25}}<\frac12.
                                                                    \label{eq:m8-globalnumeric}
\end{equation}
This proves the stated global onset \(N\ge2^{78}\).
The margin is
\(2620327/234881024>1/100\).
In every eligible stratum \(r(N)\ge3N\), and each ordered value
pair has at most \((2M(N))^2\le144N^{1/3}\) witnesses.
The count is therefore at least \(N^{2/3}/4800\).

Two additional uniform families follow from the same inequality.
For inert-supported \(m\), use \(P=2m^{1/4}\) and
\(\tau(m)\le7m^{1/6}\).
The error after \(B\) is at most
\[
                    \frac{19}{2}m^{-1/12}
                                 +\frac{69}{2}m^{-1/3}.
\]
For \(m\ge2^{60}\), its total with \(B\) is bounded by
\begin{equation}
                  \frac3{16}+\frac{19}{64}
                                      +\frac{69}{2^{21}}<\frac12.
                                                                    \label{eq:m8-coretail}
\end{equation}
Thus all \(2^em\), \(0\le e\le3\), in this family lie in \(P+P\).
At \(m\) the two values must be in different parity classes;
at \(2m\) both are in \(C\); at \(8m\) both are in \(2C\),
and division by two gives \(4m\in C+C\).
This proves the three core conclusions above the onset.

For an odd prime \(n\), use \(E(n)\le2n^{-1/2}\),
\(M(n)\le2n^{1/4}\), \(L_0=4\), and \(P=2n^{3/8}\).
The error after \(B\) is at most \(2n^{-1/8}+7n^{-1/2}\).
For \(n\ge2^{22}\), it is dominated by
\(3/10+7/2048\), and
\begin{equation}
                         3/16+3/10+7/2048<1/2.               \label{eq:m8-prime}
\end{equation}
Every such prime lies in \(P+P\); since it is odd,
it lies in \(C+2C\). The finite mixed-core coverage below will
supply all smaller odd primes, so no separate prime-family
search is required.

\subsection{Sharper selected shells}

The improved lattice bounds are
\begin{equation}
\begin{aligned}
 S_Q(X)&:=\#\{(x,y):0<Q(x,y)\le X\}\le\tfrac83 X,\\
 S_o(X)&:=\#\{(x,y):Q(x,y)\le X,\ Q(x,y)\ {\rm odd}\}\le2X.
                                                               \end{aligned}\label{eq:m8-pointsharp}
\end{equation}
For the first, use the preceding row bound and
\(\pi<22/7\), \(7/5<\sqrt2<10/7\).
At integer \(X\ge81\) the ratio is at most
\(110/49+8/21=386/147<8/3\).
The complete point table checks \(1\le X\le80\).
For odd norm the first coordinate is odd; row integration gives
\[
 S_o(X)\le\frac{\pi X}{2\sqrt2}+(1+\sqrt2)\sqrt X+1 .
\]
At \(X\ge25\), divide by \(X\) and use
\(55/49+17/35+1/25<2\).
The point table checks \(1\le X\le24\).
Both counts vanish below one, and flooring proves the real bounds.

For odd \(m\), sieve the following selected subsets:
\begin{equation}
\begin{array}{c|c|c|c|c}
 N&\text{selected subset}&\text{proportion at least}
                                      &\mathcal L&\delta\\ \hline
 m&G_2(m)&1/2-E(m)/2&8/3&1/2\\
 2m&\text{both field norms odd}&1/2&2&0\\
 4m&\text{both field norms odd}&2/3&4/3&0
\end{array}                                                    \label{eq:m8-selected}
\end{equation}
The even subset sizes are exact. Multiplication of both components
by \(\varpi\) bijects the full shell of norm \(N/2\) with the part
of the norm-\(N\) shell having both norms even.
The odd-odd complement has size \(r(N)-r(N/2)\), namely
\(4\sigma(m)\) at \(2m\) and \(16\sigma(m)\) at \(4m\).
All selected components are nonzero.

For the first row, \eqref{eq:m8-pointsharp} and
\(r_Q(t)\le2M(N)\) bound the large-prime count by
\((32/3)M(N)N/p^2\), giving normalized coefficient
\(\mathcal L=8/3\).
For the other rows use \(S_o(N/p^2)\le2N/p^2\).
The numerator is \(8M(N)N/p^2\); division by the respective
full shells \(8\sigma(m)\) and \(24\sigma(m)\) gives \(2\) and
\(4/3\). The small-prime bound \eqref{eq:m8-small} is applied
in the full shell, not in the selected subset.
Consequently, with \(\lambda\) the third-column constant
\(1/2,1/2,2/3\), a sufficient condition in each row is
\begin{equation}
       B+E(m)(\Pi(P)+\delta)+\mathcal L M(N)T(P)<\lambda,
       \quad T(P)=\sum_{\substack{p\in\mathcal I\\p>P}}p^{-2}.
                                                               \label{eq:m8-selectedtest}
\end{equation}
This follows by subtracting the union of all bad events from the
selected subset; no Hecke stability of that subset is needed.
After sieving the rows give respectively \(C+2C,C+C,C+C\).

\subsection{The complete interval and finite certificates}

For inert-supported odd \(m\), a sharper bound is
\begin{equation}
                         E(m)<\frac{18}{5}m^{-1/3}.            \label{eq:m8-Esharp}
\end{equation}
Raise the normalized prime-power factors to the sixth power:
they are
\((a+1)^6p^{5a}/\sigma(p^a)^6\).
Primes at least \(64\) contribute at most one.
For each smaller inert prime, exponents through three suffice:
the subsequent ratio is strictly less than
\((5/4)^6/p\le(5/4)^6/5<1\).
The exact product of the finite maxima is below \((18/5)^6\).
Thus \eqref{eq:m8-Esharp} is an unrestricted bound.

For the bounded \(M\)-envelope, enumerate nonincreasing exponents
on initial segments of the smallest split primes.
Swapping exponents \(a<b\) at \(p<q\) lowers the integer cost
by \((p/q)^{b-a}\), preserving the divisor product, so no
maximizer is lost. The recursion through \(2^{63}\) has \(2970\)
states. The product of the first thirteen split primes exceeds
\(2^{63}\), and \(3^{40}>2^{63}\); these bound both its depth
and exponent range.

Set \(R_0=65536\), \(D_0=2^{64}\). For \(P<R_0\), define
\begin{equation}
 V(P)=\frac1{D_0}
       \sum_{\substack{p\in\mathcal I\\P<p\le R_0}}
                    \left\lceil D_0/p^2\right\rceil
                   +\frac1{4R_0}+\frac2{R_0^2}\ge T(P).       \label{eq:m8-V}
\end{equation}
The prime list is a complete Eratosthenes sieve; the final two
terms majorize all larger primes by the two progressions.
Thus \(\Pi(P)\) is exact and \(V(P)\) is an upward rational bound.

For \(N=fm\), \(f=1,2,4\), and \(2^k\le m<2^{k+1}\), put
\[
 d_{f,k}=M(f(2^{k+1}-1)),\qquad
 e_k=\frac{18}{5\lfloor(2^k)^{1/3}\rfloor},\qquad
 P=\min\left(R_0-1,\
          \left\lfloor\sqrt{\mathcal Ld_{f,k}/e_k}\right\rfloor\right).
\]
All of these quantities are obtained by integer and rational
operations, using squared comparisons for square roots.
The certificate verifies
\begin{equation}
 \frac3{16}+e_k(\Pi(P)+\delta)
                      +\mathcal Ld_{f,k}V(P)
                            <\lambda-\frac1{256}              \label{eq:m8-intervalcert}
\end{equation}
on every interval in the following table:
\[
\begin{array}{c|c|c|c}
 f&\text{core family}&k\text{-range}&\text{number}\\\hline
 1&C+2C&27\le k\le59&33\\
 2&C+C\text{ at }2m&26\le k\le59&34\\
 4&C+C\text{ at }4m&20\le k\le59&40
\end{array}
\]
There are \(107\) tests, with all exact inputs and fractions in
the submitted certificate. Each row has minimum margin greater than
\(1/256\); Appendix~\ref{app:verifier} records the three exact minima.
Together with \eqref{eq:m8-coretail}, these tests cover all
inert cores above the finite ranges below.

It remains to describe that finite coverage independently of
any asymptotic statement. Put \(B_0=2^{27}\).
A complete prime sieve constructs the membership flags of
\(C\cap[1,B_0]\) by removing multiples of \(2\) and every
prime \(5,7\pmod8\). No larger prime can divide a value in this
interval. These complete flags are also authenticated by
primitive coordinate enumeration:
\(x\ge1\) odd, \(y\ge0\),
\(x^2+2y^2\le B_0\), \(\gcd(x,y)=1\).
Independent sign changes justify those coordinate restrictions.
The two full bytearrays are compared, rather than merely their
hashes.

Pack the flags as \(B_C=\sum_{n\in C,n\le B_0}2^n\), and form
\(B_{2C}\) by doubling their positions. The exact truncated unions
\[
 U_e=\bigvee_{\substack{a\le5000\\a\in C}}(B_C\ll a),
\]
\[
 U_o=\bigvee_{\substack{a\le5000\\a\in C}}(B_{2C}\ll a)
         \ \vee\
          \bigvee_{\substack{b\le5000\\b\in2C}}(B_C\ll b)
\]
have only genuine sums. There are \(699\) core anchors and
\(365\) doubled anchors.
After truncation at \(B_0\), \(U_e\) equals exactly the periodic
mask \(2,4,6\pmod8\); \(U_o\) equals exactly the odd mask with
position one removed. Thus every one of \(50331648\) even
strong targets and \(67108863\) mixed odd targets is covered,
with no locally forbidden position covered. A short anchor
has not been postulated: the union actually covers every
required bit. The certificate programs reconstruct these arrays
and compare the complete target masks.

\subsection{Completion by descent}

Suppose that a least eligible exception to \(C+C\) exists.
Its dyadic valuation is one or two. A split prime factor \(q\)
can be removed: the quotient is a smaller eligible even target
at least two, and multiplication by \(q\) lifts a representation.
Thus the least exception is \(2m\) or \(4m\) with odd
inert-supported \(m\). The finite coverage, the corresponding
interval row in \eqref{eq:m8-intervalcert}, and
\eqref{eq:m8-coretail} cover every such \(m\), a contradiction.

For \(C+2C\), a least missing odd \(N\ge3\) with a split prime
factor either is a single split prime, or descends to a smaller
odd target at least three. The prime boundary is covered by
\eqref{eq:m8-prime} above \(2^{22}\) and by the finite mixed
coverage below that bound. With no split factor, \(N\) is
inert-supported and is covered by the finite computation, its
\(33\) interval tests, and the infinite tail.
This proves both strong core identities. Their local exclusions
and \eqref{eq:m8-union} prove the full primitive identity.
The explicit count was already obtained from
\eqref{eq:m8-globalnumeric}, completing
Theorem~\ref{thm:minus8-exact}.
